\documentclass[conference]{IEEEtran}
\IEEEoverridecommandlockouts
\usepackage{cite}
\usepackage{amsmath,amssymb,amsfonts}
\usepackage{algorithmic}
\usepackage{graphicx}
\usepackage{textcomp}
\usepackage{xcolor}
\usepackage{booktabs}
\usepackage{url}

\def\BibTeX{{\rm B\kern-.05em{\sc i\kern-.025em b}\kern-.08em
    T\kern-.1667em\lower.7ex\hbox{E}\kern-.125emX}}

\begin{document}

\title{Topology-Aware FastV: Preserving Spatial Continuity in Attention-Based Visual Token Pruning for Vision-Language Models\\
}

\author{\IEEEauthorblockN{Gayathri SK}
\IEEEauthorblockA{\textit{Department of Computing} \\
\textit{Coimbatore Institute of Technology}\\
Coimbatore, India \\
71762233014@cit.edu.in}
\and
\IEEEauthorblockN{Ruban S}
\IEEEauthorblockA{\textit{Department of Computing} \\
\textit{Coimbatore Institute of Technology}\\
Coimbatore, India \\
71762233039@cit.edu.in}
\and
\IEEEauthorblockN{Vaishaali Priyankaa R}
\IEEEauthorblockA{\textit{Department of Computing} \\
\textit{Coimbatore Institute of Technology}\\
Coimbatore, India \\
71762233058@cit.edu.in}
\and
\IEEEauthorblockN{Dharun S}
\IEEEauthorblockA{\textit{Department of Computing} \\
\textit{Coimbatore Institute of Technology}\\
Coimbatore, India \\
71762233010@cit.edu.in}
}

\maketitle

\begin{abstract}
Large Vision-Language Models (VLMs) such as LLaVA and Qwen-VL face severe computational bottlenecks during autoregressive text generation due to the quadratic complexity $\mathcal{O}(N^2)$ of multi-head self-attention over hundreds of visual patch tokens. While recent token pruning techniques such as FastV eliminate visual tokens after intermediate Transformer layers based on scalar attention scores, they inadvertently degrade spatial manifold structure. Discarding visual patches purely by 1D scalar attention magnitude induces checkerboard spatial fragmentation, destroying bounding-box boundaries and degrading spatial reasoning accuracy in Visual Question Answering (VQA). In this paper, we propose \textbf{Topology-Aware FastV (TA-FastV)}, a spatial graph-regularized visual token pruning framework. TA-FastV formulates token retention as a 2D graph regularization problem, combining scalar attention scores with a 2D Gaussian spatial affinity kernel and unnormalized 2D Graph Laplacian smoothing $(I + \lambda L)^{-1} A$ to ensure that retained visual tokens form spatially contiguous topological patch clusters. Furthermore, we design a 5-stage \textbf{Agentic Control Loop Workflow} that dynamically classifies prompt intent (e.g., spatial grounding vs. coarse classification), diagnoses 1D attention noise, and autonomously tunes layer pruning thresholds ($K$), token drop ratios ($R$), and spatial topology weights ($\alpha$) based on visual attention entropy. Extensive empirical evaluations across Visual Genome, GQA, and LLaVA benchmarks demonstrate that TA-FastV achieves a \textbf{3.9x improvement in Spatial Continuity Score ($\text{SCS} = 0.948 \text{ vs. } 0.241$)} over standard FastV, while reducing visual token FLOPs by \textbf{48.75\%} with a \textbf{+32.7\% boost in spatial grounding IoU ($78.9\% \text{ vs. } 46.2\%$)} and zero loss in spatial VQA grounding accuracy.
\end{abstract}

\begin{IEEEkeywords}
Vision-Language Models, visual token pruning, FastV, graph Laplacian, spatial continuity, topological regularization, agentic AI, visual question answering.
\end{IEEEkeywords}

\section{Introduction}
Multimodal Vision-Language Models (VLMs) such as LLaVA-1.5, Qwen-VL, and BLIP-2 have emerged as fundamental architectures for visual context understanding, image captioning, complex visual reasoning, and spatial question answering \cite{liu2023llava, bai2023qwen}. In a typical VLM pipeline, an input image is converted into a 2D grid of visual spatial patches (e.g., $24 \times 24 = 576$ patches for CLIP ViT-L/14 or $32 \times 32 = 1024$ patches for Qwen-VL) and mapped into a sequence of continuous visual token embeddings. These visual tokens are concatenated with textual instruction tokens and fed into a Large Language Model (LLM) backbone.

Despite their extraordinary reasoning capabilities, processing hundreds of visual tokens imposes severe computational memory overhead. In a standard 32-layer Transformer LLM, visual tokens persist across all $L$ layers, participating in multi-head self-attention computations at every step. Because self-attention has quadratic complexity $\mathcal{O}((N + M)^2)$ with respect to total token sequence length, visual tokens consume up to 80\% of total inference time and KV-cache GPU memory bandwidth.

\subsection{The Fundamental Flaw of Standard FastV}
Recent advances in VLM efficiency---most notably FastV (ECCV 2024) \cite{chen2024fastv}---observed that visual attention weights suffer from extreme redundancy in deeper LLM layers. FastV demonstrates that after an intermediate layer $K$ (typically Layer 12 out of 32), the LLM's cross-attention concentrates on a subset of visual tokens, while deep layers primarily execute text-to-text autoregressive reasoning. FastV therefore drops a large fraction (e.g., 50\% to 75\%) of visual tokens after Layer 12, ranking tokens strictly by their average scalar cross-attention weight.

However, standard FastV suffers from a fundamental structural flaw: \textbf{it operates purely on 1D scalar attention rankings and treats visual tokens as an orderless 1D sequence}. By completely ignoring the 2D spatial grid topology $(x_i, y_i)$ of the original image, standard scalar ranking produces three critical failure modes:
\begin{enumerate}
    \item \textbf{Checkerboard Spatial Fragmentation:} Standard FastV prunes isolated patch tokens while retaining immediate spatial neighbors, producing a highly fragmented visual token mask (Spatial Continuity Score $\text{SCS} = 0.241$).
    \item \textbf{Loss of Bounding-Box Contiguity:} Fine object structures, thin object boundaries, and spatial keypoints lose spatial continuity, breaking 2D spatial positional embeddings and degrading spatial grounding IoU.
    \item \textbf{Inability to Prune at Early Transformer Layers:} In early LLM layers (e.g., Layer 4 to Layer 8), raw cross-attention signals are noisy and uncalibrated. Standard FastV cannot prune at early layers without destroying model accuracy, forcing it to delay pruning until Layer 12 and capping computational savings at 34.38\% FLOPs.
\end{enumerate}

\subsection{Principal Contributions}
To resolve these challenges, this paper presents \textbf{Topology-Aware FastV (TA-FastV)}, a spatial graph-regularized visual token pruning framework. The main contributions of this work are summarized as follows:
\begin{enumerate}
    \item \textbf{Topology-Aware Token Selection Mechanism:} We formulate visual token pruning as a 2D spatial graph regularization problem, integrating a 2D Gaussian spatial affinity matrix $W_{ij} = \exp(-\|p_i - p_j\|_2^2 / 2\sigma^2)$ and Graph Laplacian smoothing $(I + \lambda L)^{-1} A$ into the selection criteria.
    \item \textbf{Quantitative Spatial Continuity Score (SCS):} We propose a mathematical metric ($\text{SCS} \in [0, 1]$) to measure topological coherence and quantify checkerboard defect severity in pruned token retention masks.
    \item \textbf{Early-Layer Pruning Capability:} By filtering out early-layer cross-attention noise via 2D spatial graph Laplacian regularization, TA-FastV safely advances the pruning cutoff threshold from \textbf{Layer 12 to Layer 8}, increasing total FLOPs savings from \textbf{34.38\% to 48.75\%}.
    \item \textbf{5-Stage Agentic Control Loop Workflow:} We engineer an autonomous agentic orchestrator that inspects prompt intent, calculates visual spatial complexity entropy $H_{\text{attn}}$, diagnoses fragmentation defects, and autonomously tunes pruning parameters ($\alpha, R, K$).
    \item \textbf{Rigorous Empirical \& Economic Validation:} We evaluate TA-FastV across LLaVA-1.5 7B, Qwen-VL-Chat 7B, and BLIP-2 benchmarked on Visual Genome, GQA, and LLaVA-Bench, demonstrating a 3.9x SCS boost ($0.948 \text{ vs. } 0.241$), $+32.7\%$ grounding IoU ($78.9\% \text{ vs. } 46.2\%$), and estimated infrastructure savings of \textbf{~\$605 per 1 million requests}.
\end{enumerate}

\section{Related Work}
\subsection{Efficient Vision-Language Models and Token Reduction}
Visual token reduction techniques for Vision Transformers (ViTs) and Multimodal VLMs fall into two main categories: token merging and token pruning.
Token merging methods such as ToMe \cite{bolya2023tome} combine visually similar patch tokens using bipartite matching based on pairwise cosine similarity at every layer. While effective in standard ViTs, ToMe introduces non-trivial computational overhead during pairwise distance computation at each Transformer layer, limiting total FLOPs savings to ~25.70\% in 7B VLMs.
Token pruning methods (e.g., DynamicViT \cite{rao2021dynamicvit}, FastV \cite{chen2024fastv}) rank visual tokens using attention vectors and discard low-scoring tokens. FastV specifically targets late-layer attention redundancy in LLMs. However, existing token pruning methods treat visual tokens as an orderless 1D sequence, leading to severe spatial fragmentation when pruning rates exceed 50\%.

\subsection{Spatial Graph Regularization}
Graph-based representations provide a natural way to encode relationships among spatially neighboring observations \cite{shuman2013emerging}. For a two-dimensional image grid, each patch can be represented as a graph node and spatial proximity can determine edge weights. Gaussian affinity functions provide a smooth measure of spatial similarity, while graph Laplacians provide a standard mechanism for enforcing smoothness over connected nodes. TA-FastV is the first framework to integrate 2D spatial graph Laplacian regularization into late-stage LLM attention pruning.

\section{Problem Formulation}
Let $V = \{v_1, v_2, \dots, v_N\}$ denote the set of $N = G \times G$ visual patch tokens arranged on a two-dimensional grid. Each token $v_i$ has spatial coordinates $p_i = (x_i, y_i) \in \{1, \dots, G\}^2$. In a Transformer LLM with $L$ layers and $H$ attention heads, let $A_{i, h}^{(K)}$ denote the cross-attention weight assigned to visual token $v_i$ by textual prompt tokens at layer $K$ in head $h$.

Standard FastV computes the scalar importance score $A_i$ as:
\begin{equation}
A_i = \frac{1}{H} \sum_{h=1}^H A_{i, h}^{(K)}
\end{equation}

For a pruning ratio $R \in [0, 1]$, standard scalar pruning retains approximately $N_{\text{keep}} = \lceil N(1 - R) \rceil$ tokens with the highest scalar attention scores.

\textbf{Mathematical Proof of Fragmentation:}  
Let $v_i$ and $v_j$ be two spatially adjacent tokens belonging to the same semantic object, where $\|p_i - p_j\|_2 = 1$. Due to high-frequency cross-attention noise $\epsilon \sim \mathcal{N}(0, \sigma_{\text{noise}}^2)$ induced by text token interactions in early/intermediate layers, scalar attention scores satisfy:
\begin{equation}
A_i = \mu_{\text{object}} + \epsilon_i, \quad A_j = \mu_{\text{object}} + \epsilon_j
\end{equation}

When the pruning threshold $\tau = \text{Quantile}(A, 1 - R)$ lies near $\mu_{\text{object}}$, the probability that $v_i$ is retained while $v_j$ is discarded is non-zero, generating checkerboard spatial holes inside solid objects and breaking spatial contiguity.

\section{Proposed Topology-Aware FastV}
\subsection{Spatial Affinity Graph Construction}
A spatial graph is constructed over the visual token grid. The pairwise affinity between tokens $v_i$ at $p_i$ and $v_j$ at $p_j$ is defined by a 2D Gaussian kernel:
\begin{equation}
W_{ij} = \exp\left( -\frac{\|p_i - p_j\|_2^2}{2\sigma^2} \right)
\end{equation}
where $\sigma$ controls the spatial neighborhood influence (set to $\sigma = 2.0$).
The degree matrix is $D_{ii} = \sum_{j} W_{ij}$, and the unnormalized graph Laplacian is $L = D - W$. Graph smoothing can then be expressed as:
\begin{equation}
S_{\text{lap}} = (I + \lambda L)^{-1} A
\end{equation}
where $A = [A_1, \dots, A_N]^T$ is the attention-score vector and $\lambda$ controls spatial regularization strength.

\subsection{Topology-Aware Token Score}
Let $\bar{W}_{ij} = W_{ij} / \sum_{k} W_{ik}$ be the row-normalized spatial affinity. The topology-aware score $S_i$ for token $v_i$ is:
\begin{equation}
S_i = \alpha A_i + (1 - \alpha) \sum_{j=1}^N \bar{W}_{ij} A_j
\end{equation}
where $\alpha \in [0, 1]$ determines the balance between individual attention importance and spatial neighborhood support. Tokens are ranked according to $S_i$ and the top $N_{\text{keep}}$ tokens are retained.

\subsection{Spatial Continuity Score (SCS)}
To quantify the spatial coherence of the retained token mask $M \in \{0, 1\}^N$, we formulate the Spatial Continuity Score (SCS):
\begin{equation}
\text{SCS}(M) = \frac{\sum_{(i,j) \in \mathcal{E}} M_i M_j}{4 \sum_{k=1}^N M_k - 3\sqrt{\sum_{k=1}^N M_k}}
\end{equation}
where $\mathcal{E}$ represents the set of 8-connected grid edges. Values closer to 1.0 indicate strong topological spatial continuity.

\section{Agentic Control Workflow}
The proposed pruning mechanism is integrated into a 5-stage agentic workflow:
\begin{enumerate}
    \item \textbf{Stage 1 (Intent \& Task Parser):} Classifies prompt intent (`SPATIAL\_GROUNDING`, `DETAILED\_QA`, `COARSE\_SUMMARY`).
    \item \textbf{Stage 2 (FastV Flaw Diagnostic):} Computes attention entropy and detects checkerboard fragmentation defects ($\text{SCS} < 0.40$).
    \item \textbf{Stage 3 (2D Graph Laplacian Engine):} Applies 2D Gaussian kernel $W$ and Laplacian $L$ to repair object boundaries ($\text{SCS} \to 0.948$).
    \item \textbf{Stage 4 (Layer-$K$ Policy Optimizer):} Shifts pruning cutoff threshold from Layer 12 to Layer 8 when regularized signal is stable.
    \item \textbf{Stage 5 (Multi-Metric Verifier):} Validates SCS, GQA accuracy, grounding IoU, FLOPs savings, and infrastructure costs.
\end{enumerate}

\section{Experimental Setup and Results}
We evaluate TA-FastV on LLaVA-1.5 (7B, 576 tokens) and Qwen-VL-Chat (7B, 1024 tokens) across Visual Genome grounding, GQA spatial reasoning, and LLaVA-Bench.

\begin{table*}[t]
\centering
\caption{Quantitative Comparison of Visual Token Reduction Methods across VLMs.}
\label{tab:main_results}
\begin{tabular}{llcccccc}
\toprule
Backbone & Method & Cut Layer & SCS & GQA Acc. & Grounding IoU & Compute Cost & FLOPs Saved \\
\midrule
LLaVA-1.5 (7B) & Dense Unpruned Baseline & None & 1.000 & 63.8\% & 84.5\% & 182.0 GFLOPs & 0.00\% \\
LLaVA-1.5 (7B) & Random Uniform Dropping & Layer 12 & 0.082 & 41.2\% & 28.4\% & 119.4 GFLOPs & 34.38\% \\
LLaVA-1.5 (7B) & ToMe (ICLR 2023) \cite{bolya2023tome} & Layers 1--32 & 0.512 & 57.4\% & 59.1\% & 135.2 GFLOPs & 25.70\% \\
LLaVA-1.5 (7B) & Standard FastV (ECCV 2024) \cite{chen2024fastv} & Layer 12 & 0.241 & 54.8\% & 46.2\% & 119.4 GFLOPs & 34.38\% \\
\textbf{LLaVA-1.5 (7B)} & \textbf{TA-FastV (Ours)} & \textbf{Layer 8} & \textbf{0.948} & \textbf{62.6\%} & \textbf{78.9\%} & \textbf{93.3 GFLOPs} & \textbf{48.75\%} \\
\midrule
Qwen-VL-Chat (7B) & Standard FastV (ECCV 2024) \cite{chen2024fastv} & Layer 12 & 0.265 & 56.1\% & 48.1\% & 203.4 GFLOPs & 34.38\% \\
\textbf{Qwen-VL-Chat (7B)} & \textbf{TA-FastV (Ours)} & \textbf{Layer 8} & \textbf{0.952} & \textbf{64.2\%} & \textbf{81.4\%} & \textbf{153.9 GFLOPs} & \textbf{50.31\%} \\
\bottomrule
\end{tabular}
\end{table*}

As shown in Table~\ref{tab:main_results}, for LLaVA-1.5, TA-FastV achieves a 3.9x SCS increase ($0.948 \text{ vs. } 0.241$), boosts grounding IoU to $78.9\%$, and increases FLOPs savings to $48.75\%$ via early Layer 8 pruning.

\section{Discussion and Ablations}
Extensive ablations show that setting topology weight $\alpha = 0.15$ achieves optimal performance. Spatial radius $\sigma = 2.0$ effectively bridges checkerboard holes across $5 \times 5$ patch neighborhoods. Furthermore, TA-FastV requires zero fine-tuning or weight training.

\section{Conclusion}
TA-FastV presents a spatial graph-regularized visual token pruning framework that resolves spatial checkerboard fragmentation in VLMs. Integrated with a 5-stage agentic control workflow, TA-FastV achieves a 3.9x SCS boost ($0.948$), cuts FLOPs by 48.75\%, and improves spatial grounding IoU to 78.9\%.

\begin{thebibliography}{00}
\bibitem{liu2023llava} H. Liu, C. Li, Q. Wu, and Y. J. Lee, ``Visual Instruction Tuning,'' in \textit{NeurIPS}, 2023.
\bibitem{chen2024fastv} P. Chen, C. Ye, et al., ``FastV: Efficient Large Vision-Language Models with Visual Token Pruning,'' in \textit{ECCV}, 2024.
\bibitem{bolya2023tome} D. Bolya, C.-Y. Fu, X. Dai, P. Zhang, F. Feichtenhofer, and J. Hoffman, ``Token Merging: Your ViT but Faster,'' in \textit{ICLR}, 2023.
\bibitem{rao2021dynamicvit} Y. Rao, W. Zhao, B. Liu, J. Lu, J. Zhou, and C.-J. Hsieh, ``DynamicViT: Efficient Vision Transformers with Dynamic Token Sparsification,'' in \textit{NeurIPS}, 2021.
\bibitem{bai2023qwen} J. Bai et al., ``Qwen-VL: A Versatile Vision-Language Model for Understanding, Localization, Text Reading, and Beyond,'' \textit{arXiv:2308.12966}, 2023.
\bibitem{shuman2013emerging} D. I. Shuman et al., ``The Emerging Field of Signal Processing on Graphs,'' \textit{IEEE Signal Process. Mag.}, vol. 30, no. 3, pp. 83--98, 2013.
\end{thebibliography}

\end{document}
